% 第4章 4.5.1：W^{1,p} 嵌入与维数。(a) (p,d) 平面：p<d 时嵌入 L^{p*}，p*=dp/(d-p)；p>d 时有 C^{0,1-d/p} 代表；p=d 为临界线。
% 标出 H^1（p=2）在 d=1,2,3 的位置：d=1 得 C^{0,1/2}，d=2 临界，d=3 得 L^6。
% (b) p*=dp/(d-p) 随 p 的变化（d=2,3,4），p→d 时趋于无穷（解析式绘制）。
\begin{tikzpicture}[font=\figurefont\footnotesize]
\pgfplotsset{every axis/.style={width=52mm,height=44mm,axis lines=left,
  x axis line style={draw=FigureStroke,line width=.6pt},y axis line style={draw=FigureStroke,line width=.6pt},
  tick style={draw=FigureStroke},tick label style={font=\figurefont\scriptsize,text=FigureInk},
  title style={font=\figurefont\footnotesize},
  xlabel style={at={(axis description cs:1,0)},anchor=west,font=\figurefont\scriptsize},
  ylabel style={at={(axis description cs:0,1)},anchor=south,rotate=-90,font=\figurefont\scriptsize}}}
\begin{axis}[name=A,xmin=1,xmax=6.4,ymin=1,ymax=5.4,xtick={1,2,3,4,5,6},ytick={1,2,3,4,5},xlabel={$p$},ylabel={$d$},title={(a) 积分指数与维数}]
\fill[FigureBlueFill] (axis cs:1,1) -- (axis cs:5.4,5.4) -- (axis cs:1,5.4) -- cycle;
\fill[FigureGreenFill] (axis cs:1,1) -- (axis cs:6.4,1) -- (axis cs:6.4,5.4) -- (axis cs:5.4,5.4) -- cycle;
\draw[FigureInk,line width=.8pt,dashed] (axis cs:1,1) -- (axis cs:5.4,5.4);
\node[font=\figurefont\scriptsize,text=FigureBlue,align=center] at (axis cs:2.2,4.4) {$p<d$\\$L^{p^*}$};
\node[font=\figurefont\scriptsize,text=FigureGreen,align=center] at (axis cs:4.9,2.0) {$p>d$\\Hölder代表};
\node[font=\figurefont\scriptsize,text=FigureInk,rotate=39,anchor=south] at (axis cs:3.9,3.9) {临界$p=d$};
\addplot[only marks,mark=*,mark size=1.6pt,FigureInk] coordinates {(2,1) (2,2) (2,3)};
\node[font=\figurefont\scriptsize,text=FigureInk,anchor=west] at (axis cs:2.1,1.25) {$C^{0,1/2}$};
\node[font=\figurefont\scriptsize,text=FigureInk,anchor=east] at (axis cs:1.9,3) {$L^6$};
\end{axis}
\begin{axis}[at={(A.east)},anchor=west,xshift=14mm,xmin=1,xmax=4.2,ymin=0,ymax=16,xtick={1,2,3,4},ytick={0,4,8,12,16},xlabel={$p$},ylabel={$p^*$},title={(b) $p^*=dp/(d-p)$},
  legend style={draw=none,fill=none,font=\figurefont\scriptsize,at={(0.97,0.03)},anchor=south east,cells={anchor=west}},samples=200]
\addplot[draw=FigureBlue,line width=1.1pt,domain=1:1.778] {2*x/(2-x)};
\addplot[draw=FigureGreen,line width=1.1pt,dashed,domain=1:2.526] {3*x/(3-x)};
\addplot[draw=FigureYellow,line width=1.1pt,dash dot,domain=1:3.2] {4*x/(4-x)};
\draw[FigureMuted,line width=.5pt,densely dotted] (axis cs:2,0) -- (axis cs:2,16);
\draw[FigureMuted,line width=.5pt,densely dotted] (axis cs:3,0) -- (axis cs:3,16);
\draw[FigureMuted,line width=.5pt,densely dotted] (axis cs:4,0) -- (axis cs:4,16);
\addplot[only marks,mark=*,mark size=1.4pt,FigureGreen] coordinates {(2,6)};
\node[font=\figurefont\scriptsize,text=FigureGreen,anchor=east] at (axis cs:1.95,6.6) {$(2,6)$};
\legend{$d=2$,$d=3$,$d=4$}
\end{axis}
\end{tikzpicture}
